
Nguồn: 
\href{https://daibieunhandan.vn/giot-nuoc-do-se-khong-tan-di-406217}{Báo Đại Biểu Nhân Dân}


\textbf{“Đối với chúng con, niềm say mê của bố với khoa học, những trăn trở của bố với đất nước là lời dạy vô giá. Những lời phát biểu chân thành, dũng cảm; giản dị nhưng đầy mạnh mẽ - như mong ước của bố - sẽ như một giọt nước nhỏ bé hòa vào nhiều triệu giọt nước khác của dân tộc để tạo thành dòng thác đổi mới của đất nước. Giọt nước đó sẽ không tan đi khi bố chia xa...” - GS. Phan Dương Hiệu, con trai của cố GS. Phan Đình Diệu nói lời tiễn biệt bố mình.}


Hơn 10 ngày sau khi đưa tiễn bố, cũng là người thầy lớn của mình, GS Phan Dương Hiệu đã dành cho chúng tôi một cuộc trò chuyện về những giá trị tri thức - tinh thần mà bố anh đã trao truyền, để lại...

\paragraph{“Một thời, chúng ta đã đi nhanh”} \ \\
%\vspace{1cm}
\textit{- Theo đánh giá của GS.TS Hà Huy Khoái, nguyên Viện trưởng Viện Toán học Việt Nam thì trình độ tin học của Việt Nam những năm 1970 là cao hơn so với các nước Đông Nam Á khác, mà công đầu là nhờ GS Phan Đình Diệu - “người anh cả của nền tin học Việt Nam”. Tiếc rằng: “Chúng ta từng đi trước, nhưng vì nhiều nguyên nhân nên đã tụt lại phía sau”. Nguyên nhân, theo anh là do đâu?}

- Muốn phát triển một lĩnh vực, điều quan trọng nhất vẫn là con người. Tôi được nghe các chú trong Viện Khoa học Tính toán \& Điều khiển (nay là Viện CNTT Việt Nam, nơi GS Phan Đình Diệu là Viện trưởng đầu tiên - PV) kể lại, bố tôi có khả năng truyền cảm hứng và thu hút những người giỏi, trẻ về Viện. Nhờ thế, Viện đã có một đội quân chất lượng, đầy nhiệt huyết, say mê phát triển các ý tưởng mới. Đất nước bấy giờ cũng có những người tâm huyết lãnh đạo lĩnh vực giáo dục, nghiên cứu như cụ Tạ Quang Bửu, Trần Đại Nghĩa... Một khi mà người đi đầu, cộng với cơ chế cho phép việc thu hút và trọng dụng con người dựa trên thực lực tài năng và tâm huyết của họ, chúng ta đã đi nhanh. Ngược lại, khi có những rào cản cho việc trọng dụng và sử dụng tài năng, làm người giỏi phải ra đi, việc tụt hậu trong một thế giới đầy năng động, thay đổi từng giờ, là không tránh khỏi.

\textit{- Cũng theo lời GS Hà Huy Khoái: “Sự phát triển, những vấn đề mà xã hội đang phải đối mặt gần như là “minh họa” những gì GS Phan Đình Diệu từng nói từ mấy chục năm trước”. Điển hình là quốc nạn tham nhũng, mà bố anh từng nhận định: “Tham nhũng về thực chất là biến một cách phi pháp quyền lực thành hàng hóa”. Ngày hôm nay, khi chứng kiến những vấn nạn chạy chức chạy quyền, những đại án tham ô tham nhũng... theo anh, cần phải làm gì để giúp kiểm soát được “thị trường quyền lực” đó?}

- Những vấn nạn chạy chức chạy quyền, những đại án tham ô tham nhũng ngày hôm nay chứng tỏ: Quyền là một thứ hàng hóa siêu lợi nhuận và người ta phải đút lót (một cách đầu tư) để chạy quyền, rồi khi có nó thì tận thu (tham nhũng).

Chúng ta vui mừng khi những đại án được đưa ra ánh sáng. Nhưng đó chỉ là bề nổi, để hướng tới triệt tận gốc vấn nạn này thì cần sự đổi mới mạnh mẽ để tăng sự minh bạch. Và để được như thế thì cần phải bỏ sự độc quyền, không thể có minh bạch nếu vừa làm vừa tự kiểm tra. Đối với tôi, độ minh bạch nhân với độ tham nhũng có “tích không đổi”. Không có xã hội nào tuyệt đối minh bạch, do vậy cũng không có xã hội nào không có tham nhũng. Nhưng tham nhũng sẽ trầm trọng hơn ở những nơi độ minh bạch thấp.

\textit{- “Bố nói với con, cuộc sống cần nhất sự trung thực. Tưởng chừng đơn giản nhưng trung thực bao gồm cả dũng cảm, trung thực với chính mình để có những chính kiến độc lập, và trung thực trong cuộc đời để dũng cảm nói lên những ý kiến tâm huyết...” - Chứng kiến và trải nghiệm nào đã khiến anh nghĩ rằng: Để sống trung thực chưa bao giờ là điều đơn giản?}

- Không chỉ là những trải nghiệm riêng lẻ, tôi ngấm dần điều đó trong cuộc sống, qua những việc làm của bố tôi. Sống trung thực đòi hỏi tâm thế sẵn sàng chấp nhận những mất mát trong sự nghiệp và cuộc sống để có thể đối diện và lên tiếng thẳng thắn trước những vấn đề gai góc của xã hội.

\textit{- Anh từng viết: “Trong mật mã, việc xây và phá xảy ra liên tục và bổ trợ lẫn nhau để phát triển ngành”. Nhưng trong những ưu tư trăn trở của bố anh trước công cuộc đổi mới của đất nước, việc “người xây kẻ phá” lại bị coi là một rào cản kìm hãm sự phát triển. Đã bao giờ anh liên hệ thế?}

- Xây và phá trong mật mã giống như sự phản biện trước một lý thuyết mới, khác với ý “người xây kẻ phá” trong cuộc sống. Phá ở đây là để cho thấy lý thuyết anh đưa ra không thuyết phục, vì thế anh cần phát triển nó tốt hơn. Trong mật mã, có những sơ đồ đứng vững trong hàng thập kỷ rồi bị phá bởi những công cụ tấn công mới. Nếu sự phá là công khai thì mọi người đều biết để tránh và xây sơ đồ mới tốt hơn. Còn nếu phá mà bí mật thì có thể gây nhiều tác hại, ta cứ dùng cái mà ta tưởng là tốt, nhưng thực chất lại đã bị phá và do đó không còn bảo mật được thông tin.

Cũng như vậy, nhiều lý thuyết trong cuộc sống cần được phản biện. Cuộc sống thay đổi, khó có lý thuyết nào là bất di bất dịch. Nếu sự phản biện là công khai, ta sẽ biết nên giữ cái nào, bỏ cái nào để phát triển. Còn nếu sự phản biện bị hạn chế thì rất dễ khiến ta đi theo những lý thuyết lỗi thời. 
\paragraph{“Tại sao lại áp lực khi có một ông bố giỏi?”}\ \\
%\vspace{1cm}
\textit{- Chiếc máy tính đầu tiên ở Việt Nam đã được tạo ra bởi bố anh và cộng sự, với sự trợ lực từ những người bạn Pháp. Đó có phải là lý do anh chọn nước Pháp làm nơi học tập và tu nghiệp?}

- Không, lựa chọn của tôi không liên quan gì tới câu chuyện chiếc máy tính đầu tiên của Việt Nam cả. Nhưng đúng là bố tôi có nhiều kỷ niệm đẹp với nước Pháp, cũng như các cộng sự Pháp và tự một lúc nào đó tôi đã quý mến nước Pháp hơn là những nơi khác.

\textit{- Ở vào thế hệ anh, cũng như sau này, hầu hết dân IT đều coi nước Mỹ là thiên đường của CNTT. Vì sao anh vẫn gắn bó với nước Pháp?}

- Trong ngành mật mã tôi theo đuổi thì Pháp là một trung tâm có vị thế khá vững ở hàng đầu trên thế giới. Môi trường sống, bề dày văn hóa và lịch sử châu Âu rất hấp dẫn, xã hội Pháp tôn trọng quyền con người. Một khi xa quê hương, tôi thấy Pháp là nơi thích hợp cho việc phát triển khoa học và cuộc sống.

\textit{- Bảo vệ tiến sĩ ở một trường tốt nhất châu Âu và hàng đầu thế giới, từng đảm nhiệm vị trí Chủ tịch hội nghị của một trong những hội nghị lớn nhất thế giới về mật mã... hành trang anh có được là gì, và món quà nào anh có thể tặng cho đất nước?}

- Đó là những hoạt động bình thường của một người làm khoa học và nó gây hứng thú cho nhà khoa học. Tôi không bao giờ nghĩ sẽ đem lại món quà gì cho đất nước. Nếu có thể đóng góp gì đó có ích thì cũng là nhu cầu tự nhiên của một người con luôn muốn gần gũi quê hương, giúp bản thân có thêm động cơ để phấn đấu và phát triển.

\textit{- Tục ngữ Việt Nam có câu: “Con hơn cha là nhà có phúc”. Đã bao giờ anh cảm thấy bị áp lực bởi quan niệm đó và cái bóng quá lớn của bố mình?}

- Chưa bao giờ. Tại sao lại bị áp lực vì có một ông bố giỏi nhỉ? Tôi chỉ thấy may mắn và hạnh phúc vì bố tôi rất say mê công việc và yêu thương vợ con. Tôi nghĩ câu tục ngữ đó mang tính động viên hơn là gây áp lực. “Mọi lý thuyết đều có thể sai”, và câu tục ngữ đó cũng có thể sai.  

\textit{- Thay vì ẩn mình trong “tháp ngà khoa học”, bố anh dường như luôn chủ trương một thứ “khoa học phục vụ cuộc sống”.  Ở vị trí một nhà nghiên cứu, học thuật, anh nghĩ mình có thể sống một cuộc đời như bố để cùng hướng đến điều đó?}

- Khái niệm “tháp ngà khoa học” vẫn còn nhiều tranh cãi. Những kết quả lý thuyết số rất đẹp đã từng được chiêm ngưỡng như những kết quả thuần tuý lý thuyết, nay lại là cơ sở cho những ứng dụng mật mã mà ta dùng hàng ngày, vậy nó là tháp ngà hay ứng dụng? Tháp ngà của quá khứ nhưng lại đem tới ứng dụng trong tương lai.

Từ những năm 1970, bố tôi đã nhìn nhận tương lai sẽ có cuộc cách mạng về xử lý thông tin với máy vi tính và ông say mê cổ vũ cho sự phát triển tin học. Lựa chọn của ông là sự kết hợp giữa đam mê khoa học và mong muốn phát triển những ứng dụng có ích cho đất nước. Để có những ứng dụng hữu ích, một nền tảng lý thuyết vững chắc là rất cần thiết nên ông đã xây dựng đồng thời những nhóm lý thuyết và ứng dụng trong Viện Khoa học Tính toán và Điều khiển ngay từ bước khởi đầu.  
Tôi cũng rất muốn hướng những nghiên cứu của mình vừa mang tính lý thuyết sâu vừa có những ứng dụng có ý nghĩa.

\textit{- Đằng sau một sự nghiệp đồ sộ và một bầu nhiệt huyết cống hiến hết mình cho đất nước, những món quà nhỏ mà bố anh dành riêng cho vợ con là gì?}

- Đó chính là sự hài hước, giản dị, bình thản trước mọi vấn đề. Đêm nào cũng làm việc khuya nhưng bố luôn tạo cảm giác ấm cúng cho gia đình. Bố tôi rất thích đọc sách, ngâm thơ và làm thơ, giọng bố rất truyền cảm. Tôi rất nhớ những đêm cả nhà quây quần nghe bố đọc thơ....

\textit{- Cách bố anh nuôi dạy ba người con thành đạt có giúp ích nhiều cho “cuốn giáo trình làm bố” của anh? }

- Bố tôi giản dị lắm, ông không nói gì to tát với tôi nhưng có cách nói truyền cảm hứng kỳ lạ. Mỗi khi đi công tác, bố thường mua cho chúng tôi những cuốn sách và đĩa nhạc hay, cả những món đồ chơi độc đáo vừa mới có trên thế giới như rubik, máy chiếu hình ảnh để giúp khơi gợi trí tò mò. Ông luôn để con cái phát triển tự nhiên theo sở thích. Tôi nghĩ rằng chúng tôi cũng sẽ làm như ông, để con cái chủ động trong việc học nhưng tôi cũng dành nhiều thời gian để chơi, trò chuyện cùng con, cùng tham gia các buổi hòa nhạc, triển lãm, tìm hiểu khoa học...

Mặc dù lúc này người ta hay nói đến khái niệm “công dân toàn cầu”, nhưng trong thâm tâm tôi vẫn luôn nghĩ nếu một đứa trẻ có ý thức về quê hương nguồn cội thì sau này sẽ vững vàng trong cuộc sống. Vì thế, năm nào chúng tôi cũng đưa con về với ông bà hai tháng hè, đi học hè cùng các bạn Việt Nam.

Và hơn hết, sự chính trực và cách sống thanh thản lạc quan của bố đã cho chúng tôi một niềm tin vào Chân lý, vào cái Đẹp của cuộc sống. Đó cũng là điều chúng tôi mong muốn con mình có thể cảm nhận.

\textit{- Tên anh có phải do bố anh đặt? Anh nghĩ điều gì đã được gửi gắm vào cái tên đó?}

- Đúng vậy, tên bố tôi là Diệu, tên mẹ tôi là Hương, nên bố mẹ tôi lái lại là Dương Hiệu (các chị tôi cũng tên là Dương), chắc ý mong tôi... đẹp trai như bố và thông minh như mẹ (cười). 

- Xin cảm ơn những chia sẻ nhiệt thành của anh!

Lê Quân thực hiện